\section{アルゴリズム：SC-DenseAttrib}
\label{sec:algorithm}

\Cref{alg:sc-dense} に SC-DenseAttrib アルゴリズムの全体を示す．
本手続きは5つの段階から構成される：サポート系列の計算，ドナープールの構築，重みの推定，因果効果の推定，およびプラセボベースの推論である．

\begin{algorithm}[t]
\caption{SC-DenseAttrib：密集パターン帰属のための合成コントロール}
\label{alg:sc-dense}
\begin{algorithmic}[1]
\REQUIRE トランザクション系列 $\mathcal{D}$，処置パターン $P^*$，候補パターン $\mathcal{P}$，介入時刻 $t_0$，ウィンドウサイズ $W$
\ENSURE 因果効果 $\hat{\tau}$，p値 $p$
\STATE \Cref{def:windowed-support} によりサポート系列 $\mathbf{s}_{P^*}$ を計算
\STATE ドナープール $\mathcal{D}(P^*) \leftarrow \{P \in \mathcal{P} : P \cap P^* = \emptyset\}$ を構築
\FOR{各 $P_j \in \mathcal{D}(P^*)$}
    \STATE $\mathbf{s}_{P_j}$ を計算し，分散によりフィルタリング
\ENDFOR
\STATE ドナー行列 $\mathbf{D} \in \mathbb{R}^{T \times J}$ を構成
\STATE 重み $\mathbf{w}^* \leftarrow \arg\min_{\mathbf{w} \in \Delta^J} \|\mathbf{s}_{P^*}^{\mathrm{pre}} - \mathbf{D}^{\mathrm{pre}} \mathbf{w}\|^2$ を推定
\STATE 反事実 $\hat{\mathbf{s}}^{\mathrm{CF}} \leftarrow \mathbf{D} \mathbf{w}^*$ を計算
\STATE すべての $t$ に対して $\hat{\tau}(t) \leftarrow \sigma_{P^*}(t) - \hat{\sigma}^{\mathrm{CF}}(t)$ を計算
\STATE $r^* \leftarrow \mathrm{RMSPE}_{\mathrm{post}}(\hat{\tau}) / \mathrm{RMSPE}_{\mathrm{pre}}(\hat{\tau})$
\FOR{各ドナー $P_j$ をプラセボ処置ユニットとして}
    \STATE プラセボ効果と比 $r_j$ を計算
\ENDFOR
\STATE $p \leftarrow |\{j : r_j \geq r^*\}| / (J + 1)$
\RETURN $\hat{\tau}$, $p$
\end{algorithmic}
\end{algorithm}

\subsection{重みの推定}

式~\eqref{eq:weight-opt} の最適化は単体上の凸二次計画法である．
これを Frank-Wolfe（条件付き勾配）アルゴリズムで解く．各反復では，勾配計算に $O(J)$，頂点の特定に $O(J)$ を要し，重み推定全体の計算量は $O(KJ)$ となる．ここで $K$ は反復回数（通常 $K \leq 1000$）である．

\subsection{ドナープールの構築}

アイテム非共有基準（\Cref{def:donor-pool}）により，$P^*$ 内のアイテムを対象とする介入の直接的影響を受けない制御パターンが保証される．
さらに，介入前期間における最小分散（退化した平坦な軌跡を回避するため）および，オプションとして処置系列との介入前相関によりドナーをフィルタリングする．

\subsection{計算量解析}

\begin{theorem}[計算量]
\label{thm:complexity}
SC-DenseAttrib の計算量は $O(N \cdot |\mathcal{P}| \cdot |P_{\max}| + J^2 \cdot t_0 + J \cdot (J^2 \cdot t_0))$ である．ここで $|P_{\max}|$ は最大パターン長，$J = |\mathcal{D}(P^*)|$ はドナープールサイズ，$t_0$ は介入前期間の長さである．
\end{theorem}

\begin{proof}
サポート系列の計算はパターンあたり $O(N \cdot |P|)$ を要する．
ドナープール構築は集合交差チェックに $O(|\mathcal{P}| \cdot |P^*|)$ を要する．
重み推定は二次形式に $O(J^2 \cdot t_0)$ を要する．
プラセボ検定は重み推定を $J$ 回繰り返し，各回 $J$ 個のドナーを用いるため，合計 $O(J \cdot J^2 \cdot t_0)$ となる．
\end{proof}
